\chapter{Neural Network pricer for Heston vanilla options}
Nuestro objetivo en este apartado es construir un surrogate pricer para opciones vanilla bajo el modelo de Heston con volatilidad estocástica usando redes neuronales.\\

El pricing clásico, como es COS (FFT, PDE, etc) es muy preciso, pero es costoso.
Sin embargo, la red neuronal aproxima el mapa directo $(\Theta; m, \tau, r)\to \sigma_{imp}$.
Es importante notar que en este capitulo nos centramos en la red de pricing, no de calibración.\\

The objective of the network is to learn a smooth and accurate approximation of the implied volatility as a function of the model parameters and contract characteristics.

%%%%%%%%%%%%%%%%%%%%%%%%%%%%%%%%%%%%%%%%%%%%%%%%%
\section{Problem Formulation}

%%%%%%%%%%%%%%%%%%%%%%%%%%%%%%%%%%%%%%%%%%%%%%%%%
\section{Synthetic Data Formulation}
We need data to train the neural network. Para tener una fuente controlada y reproducible, vamos a trabajar con datos sintéticos, generados a partir del método de COS, para alimentar a la red neuronal.\\

Los datos que la red necesita son vectores de tamaño $n$, donde hay $n-1$ parámetros en las primeras columnas, y en la $n$-ésima tenemos la volatilidad implícita, calculada con el método de COS y Brent usando los $n-1$ parámetros.

Esto conlleva un problema principal, y es que vamos a trabajar en alta dimensionalidad: tenemos el set de los parámetros de Heston, $\Theta = \{\rho, \kappa, \nu_0, \bar\nu, \gamma\}$, y el set de los parámetros dados por la opción y el mercado, $\Sigma = \{K, S_0, \tau, r\}$. Esto quiere decir que tenemos que generar un espacio $\mathbb R^9$ de 9 dimensiones. 

Si quisieramos crear un grid "clásico" para poblar este espacio, tomando $N$ puntos de cada parámetro, necesitaríamos $N^9$ puntos, Solamente con tomar $N=100$, ya se vuelve un problema extremadamente costoso.\\

Para reducir un poco este problema, si fijamos el strike price en $K=1$, podemos trabajar en unidades de \textit{moneyness}, $m=S_0/K$, De esta forma, ahora trabajamos en $\mathbb R^8$.\todos{ TODO: quizas esto ponerlo en 3.2.1.}\\

Aún así, el problema de poblar este espacio sigue siendo muy costoso. Por eso, utilizaremos un método de sampleo muy útil en este tipo de problemas, conocido como el \textit{Latin Hypercube Sampling} (LHS).

\subsection{Parameter Ranges And Constraints}

\subsection{Sampling Strategy: Latin Hypercube Sampling}
The training of the neural network requires generating a large set of model parameters distributed over a high-dimensional domain. In our case, the parameter space is $\mathbb R^8$. A purely random sampling of this space may generate highly populated regions while leaving other regions scarcely explored. As a consequence, the resulting training dataset may be poorly representative, leading to suboptimal training performance.

To avoid this problem, we will use the \textit{Latin Hypercube Sampling} (LHS), a sampling method designed to uniformly cover each dimension. The main idea consists in splitting each dimension of the parameter space into $N$ disjoint subintervals of equal length, and randomly selecting exactly one sample within each subinterval.

Thus, LHS guarantees a one-dimensional marginal distribution along each coordinate, independently of the total dimension of the space. This results in a sampling scheme that is more representative of the parameter space.

This creates a representative distribution sampling of the parameter space. In contrast to a full-grid sampling, the total number of samples $N$ remains fixed, preventing exponential growth in the number of dimensions.

\todos{poner referencia?}


\subsection{Dataset Splitting And Preprocessing}

%%%%%%%%%%%%%%%%%%%%%%%%%%%%%%%%%%%%%%%%%%%%%%%%%
\section{Neural Network Architecture}

\subsection{Model Class}
Following \cite{liu}, the pricer is implemented as a fully connected feed-forward neural network (FCNN). This class of models is particularly well suited for option pricing problems where the input consists of a low-dimensional vector of continuous variables and no explicit spatial or temporal structure is present.\\

The neural network approximates the nonlinear mapping $$f_{NN}:\mathbb R^8\to \mathbb R,$$ where the input vector is given by $$\mathbf x = (\rho, \kappa, \gamma, \bar\nu, \nu_0, m, \tau, r),$$ and the scalar output corresponds to the implied volatility of a European vanilla option.\\

The network is used as a surrogate pricer, replacing the numerical evaluation of the Heston model combined with the implied volatility inversion.

\subsection{Architecture Details}

\subsection{Output Interpretation}

%%%%%%%%%%%%%%%%%%%%%%%%%%%%%%%%%%%%%%%%%%%%%%%%%
\section{Training Procedure}
\todos{TODO: rellenar}

\subsection{Loss Function}

\subsection{Optimizer And Hyperparameters}

\subsection{Regularization And Callbacks}